\section{Principles}\label{sec:principles}

This section collects what the examples have in common: the shape that every proof takes, the rules
that the matcher imposes on a lemma written for a program, the discipline that keeps such lemmas sound and
the rewriting terminating, the diagnostics that located each obstacle, and what the same principles look
like on the C programs of the MonoCera benchmarks (\texttt{benchmarks/monocera}). The sources are the numbered
findings of \texttt{NOTES.md}, the experiment threads kept in \texttt{pending/} (preludes, programs with wrong
postconditions, scripts and logs, Oct 10 and 11, 2026), \texttt{doc/teq.md}, and the summary and case studies
of the MonoCera tests. The state of the folder on Oct 12, 2026: 27 of its 30 files verify with
\texttt{-teq -simplify} (nine of them with \texttt{-include teq:nested}); open are \texttt{is\_sorted.dfy}
(one substitution of the exit into an argument of $\sorted$), the existential postcondition of
\texttt{max\_element.dfy} (no witness), and the lower bound (Section~\ref{sec:lb}, blocked by the engine).

\subsection{The shape of a proof}\label{sec:principles:shape}

Every verified example is closed by lemmas of the following kinds, in this order from generic to specific.
The first two kinds are the database's, the others are the algorithm's knowledge, stated at the beginning of
its file.
\begin{enumerate}[nosep]
\item \textbf{The database resolves single loops.} Counters, sums, stores, selects, counts, minima and
  maxima are introduced and accelerated by the modules of \texttt{-teq} without anything written for the
  program (the four loops of the counting sort up to the table of start positions, Section~\ref{sec:cs}).
\item \textbf{Fragment lemmas resolve inner counters.} Theorem~4.8 of the paper turns the trace equation
  system of an inner loop into fresh sequences inside the iteration constraint of the outer loop, which the
  simplifier distributes into fragments over incomparable dependency lists that do not see each other's
  facts (the two minimal problems, \texttt{invs\_counter\_equation.dfy} and \texttt{invs\_fragments.dfy}). The
  module \texttt{nested} does the resolution by rewriting, one fragment at a time in the order in which
  \texttt{-teq} produces them: the inner trace is a counter, its final value is its length, the test bounds
  the length, the exit fixes it, the numbers of inner iterations are a constant sequence
  (Sections~\ref{sec:nsum} and~\ref{sec:sel:frag}). The same lemmas serve every nested example.
\item \textbf{A summary states what an inner loop computes.} Its left-hand side is the outer trace, read
  through the current value, or the whole atom of the postcondition; its side conditions are the fragments,
  matched against the facts; its right-hand side is a trace of an abstract step, $\argmin$ swaps,
  $\mathsf{insert}$, $\mathsf{shake}$, $\neg\mathsf{nodup}$ (Sections~\ref{sec:sel:sum}, \ref{sec:cock},
  \ref{sec:fd}). An inner result stored at the outer position needs no summary but a chain of specific
  substitutions (Section~\ref{sec:rows}).
\item \textbf{An induction lemma states the invariant of the outer loop} as a quantified side condition over
  the abstract step, which the simplifier decides by step lemmas on atoms; the invariant is a predicate
  indexed by the iteration ($\mathsf{bubbled}$, $\mathsf{shaken}$, the sorted prefix), and the lemma
  constrains all steps of the trace (Sections~\ref{sec:sel:ind} and~\ref{sec:bs}). This is the scheme of
  \lemma{current\_trace\_bounded} of the database: no invariant is inferred, its shape is the lemma's
  knowledge, and the database proves the step.
\item \textbf{Rules on the atoms of the postcondition} characterize a named result: the bounds of
  $\argmin$ and $\mathsf{lb}$, the regions of a partition, the counts of $\mathsf{merged}$, the two cases of
  the exit for $\mathsf{distincts}$. The postcondition is never rewritten as a whole; each atom finds its
  rule, with the exit of the loop among the facts.
\item \textbf{Name the normal form.} The counter of a loop is not inlined into the facts, so a count over
  $\bl{I}$ and over $\bl{\cnt(0, 1, n)}$ differ (\texttt{count\_occurrences.dfy}); a lemma that names the
  term makes both sides equal. Quantified preconditions are stated as atoms of the specification
  ($\mathsf{keys\_below}$, Section~\ref{sec:cs}), postconditions in the database's vocabulary
  ($\mathsf{occurs}$, $\mathsf{distincts}$, $\mathsf{count}$ over $\mathsf{selects}$) where the standard
  form has no witness or no matchable shape.
\end{enumerate}
A data-dependent number of iterations (merge, gnome sort, the early exit of the first duplicate) is no
obstacle to this shape as long as the property can be stated on the final values and the exit is read from
the facts; what the paper's loop shapes exclude is the characterization of the traces, not the postcondition
(Section~\ref{sec:gnome}).

\subsection{Writing a lemma that the matcher accepts}\label{sec:principles:matcher}

The matcher of the paper (Section~5.2) relates the binding structure of combinators, follows the definitions
of the context on demand, synthesizes dependencies that are variables of the rule as maps, and matches side
conditions against the facts while they have an open variable. From these four mechanisms follow the
patterns below, each found by an example that did not close without it (Table~\ref{tab:principles}).

\begin{table}[t]
\centering\small
\begin{tabular}{@{}lll@{}}
\toprule
pattern & found on & finding \\
\midrule
specific left-hand side of an $\mathsf{all}$ rule & cocktail sort, nested mult & 15, 29 \\
fresh variables for every $\mathsf{all}$ side condition & bubble sort, cocktail sort & 9, 27 \\
open trace variables for the traces of the goal & merge, counting sort, first duplicate & 14, 31 \\
the pin for literal patterns & merge, insertion sort, gnome sort & 14 \\
the whole atom as left-hand side & merge, gnome sort & 16 \\
arithmetic in the goal's form & cocktail sort, bubble sort & 11, 28 \\
side conditions in the forms of the included modules & insertion sorts, cocktail sort & 22, 37 \\
no condition on the rule's own term & lower bound, bubble sort & 20, 32 \\
the exit of an outer loop as two implications & first duplicate & 31 \\
atoms for constraints and quantified hypotheses & linear search, counting sort & 13, 19 \\
variables of the method named apart & fragments, invs\_fragments & 9, 21 \\
\bottomrule
\end{tabular}
\caption{The patterns of Section~\ref{sec:principles:matcher}, the examples that forced them, and the
findings of \texttt{NOTES.md} that record them.}\label{tab:principles}
\end{table}

\begin{enumerate}[nosep]
\item \textbf{A specific left-hand side.} An $\mathsf{all}$ pattern whose body has a bare bound variable of
  a variable dependency on one side, $n \le j$ or $j = \intf(nn)$, matches every comparison, respectively
  equation, fragment: the dependency is synthesized as a map of the fact's dependencies. With an
  $\mathsf{all}$ side condition such a rule cycles (Section~\ref{sec:principles:termination}). A function
  symbol on the side of the bound variable, $\intf(nn) \le n$ or $J = \cnt(0, 1, nn)$, an $\mathsf{all}$ with
  a specific body, or an atom of the specification is specific.
\item \textbf{The whole atom as left-hand side.} A variable argument of an uninterpreted function is not
  expanded by its definition when the simplifier rewrites below it: $\sorted(c, 0, n)$ with $c = \cur(C)$
  is not reached by a rule on $\cur(\tr{\ldots}{\ldots}{\ldots})$, but
  $\sorted(\cur(\tr{\ldots}{\ldots}{\ldots}), 0, s)$ is, and the length is a variable $s$ with $s = m + n$
  as a side condition, since a pattern $m + n$ of two variables is not matched (Section~\ref{sec:merge}).
\item \textbf{Fresh variables for every $\mathsf{all}$ side condition.} A condition is matched against the
  facts, with backtracking, while a variable of it is open; once all are bound it is closed and decided by
  the simplifier, which does not find a closed $\mathsf{all}$ among the facts (these carry unused
  dependencies, which the simplifier also drops from the instance, and their bound variables are renamed).
  So each $\mathsf{all}$ condition binds a fresh sequence, $\al{ks, js_6, Js_6}{k, j, J \mapsto j = \cur(J)}$,
  and an equation $js_6 = js_4 \wedge Js_6 = Js_3$ identifies it afterwards; the equation is closed and
  trivially decided. The 21st condition of the cocktail summary stayed closed for a day
  (Section~\ref{sec:cock}).
\item \textbf{Open trace variables.} The traces of the goal are variables $K$, $I$, $R$ with their equations
  as hypotheses. $\bl{\cnt(0, 1, m)}$ as a \emph{dependency} of a pattern matches $\bl{K}$ through the
  context; as an \emph{argument} of a function it never does, and a closed equation $is = \bl{\tr{\ldots}{\ldots}{\ldots}}$
  is decided and fails, since $\bl{I}$ is not expanded. A rule takes $Kv = \cnt(0, 1, m)$ or
  $Iv = \tr{0}{\ldots}{\ldots}$ as an \emph{open} condition, matched against the hypothesis, and uses
  $\bl{Kv}$, $\bl{Iv}$ afterwards (Sections~\ref{sec:merge}, \ref{sec:cs}, \ref{sec:fd}).
\item \textbf{The pin.} A literal pattern that the matcher does not handle, a disjunctive guard, the
  position $j - 1$ of a store, the initial value $n - k - 2$ of a counter, the step $-1$, becomes a
  higher-order variable $G(i, j)$, $P(j)$, $X(k)$, $d$, which matches anything, and a quantified side
  condition $\forall i, j.\ G(i, j) \Leftrightarrow \varphi$ or an equation pins it to the intended term; the
  simplifier decides the pin by beta reduction. A pin must mention only variables that the left-hand side
  or an earlier condition binds, and it must follow the fragments: a test $0 < j \wedge y[j] < y[j - 1]$ is
  split into two constraints, so its pin is two pins (\texttt{insertion\_sort.dfy}). A higher-order variable
  for a bare array read $a[i]$ is instantiated by the array itself (finding~10), so reads are pinned, as
  $U(i, j) = a[i]$ in the merge, or written literally.
\item \textbf{Arithmetic in the goal's form.} The decision procedure of the context combines difference
  constraints; $len \le 2\,\intf(m) + 1$ is not derived from $n \le \intf(m) + lo + 1$ and $lo = \intf(m)$.
  A condition states the bound as the goal has it, with the program's variable as an open rule variable
  matched against its equation ($l = \intf(m)$, then $len \le \intf(m) + l + 1$); a length $len - k - 1$
  that a pattern cannot contain is a variable $p$ with $p = len - k - 1$ as a closed condition, decided once
  $p$ is bound (Section~\ref{sec:bs}). Arithmetic on a bound variable is admissible on right-hand sides and
  in side conditions, never in a left-hand side.
\item \textbf{Side conditions in the forms of the included modules.} The fragments are rewritten by the
  modules before a summary sees them, so a summary is written for the normal form of the whole database
  that the run includes: $j = \cur(J)$ without \texttt{teq:nested}, $j = \intf(nn) + k + 1$ with it, as the
  goal prints them. A file includes exactly the modules its lemmas are written for; all nested loops of the
  folder are now written for \texttt{teq:nested}, so that the module can become part of \texttt{-teq}
  (finding~37).
\item \textbf{No condition on the rule's own term.} The instance of a side condition is simplified before it
  is decided or matched, so a condition $l = \cur(Lv)$ of a rule on $\cur(Lv)$ with $Lv$ bound rewrites
  $\cur(Lv)$ by the rule itself and does not terminate. The summaries take $\cur(\mathsf{results}(\ldots))$
  as left-hand side rather than $\mathsf{results}(\ldots)$ for the same reason, and the exit of a loop is
  read from the facts ($x = \cur(Iv)$ open, then $n \le x$, Section~\ref{sec:gnome}) or abstracted by an
  atom, never stated on the current value of the trace being rewritten (Section~\ref{sec:lb}).
\item \textbf{The exit of an outer loop.} The exit of an inner loop is a fragment and a side condition like
  the others. The exit of the outer loop stands in the conclusion of the goal, $\neg exit \vee post$, which
  \texttt{-split} does not split when the exit is a disjunction, and a rule on an atom of $post$ sees
  neither literal; a postcondition stated as two implications gives each goal its literal
  (Section~\ref{sec:fd}). The case of no iteration appears only after a first simplification, hence
  \texttt{-simplify -split -simplify} where it matters (\texttt{is\_sorted.dfy}, \texttt{reverse.dfy}).
\item \textbf{Atoms for iteration constraints and quantified hypotheses.} An iteration constraint is never a
  side condition (Section~\ref{sec:principles:termination}); a definition lemma abstracts it by an atom
  with a body that the matcher cannot synthesize, an explicit read $\mathsf{earlier}(\natf(k), js) = c$, and
  the lemma that needs the constraint takes the atom (Section~\ref{sec:lsearch}). A quantified hypothesis is
  matched by no side condition, open or closed; the specification provides the atom (Section~\ref{sec:cs}).
\item \textbf{Variables named apart.} Rule variables are not renamed apart from the variables of the goal.
  A method variable with the name of a rule variable of an open condition makes the condition coincide
  with the fact literally; it is decided true without binding the variable, and the run ends with
  ``variables of a rule not instantiated by its match'' on stderr, which \emph{replaces the goal}: zero
  goals with an error is not a proof. The fragments of \texttt{-teq} are named $K?_{14}$ and never clash;
  test methods and pins name their variables $KK$, $NN$, $len$ apart.
\item \textbf{Dependencies.} Lemmas carry the extra dependencies that the simplifier keeps in maps (the
  counter of the loop as $hs$), the matcher does not backtrack over the selection of dependencies when a
  side condition fails (hence \lemma{invs\_current\_counter\_other}), equal counters $K = I$ of one loop are
  not identified, and a higher-order summand $X(k)$ matches a single summand only, which is why the
  resolved forms are written literally.
\end{enumerate}

\subsection{Soundness}\label{sec:principles:sound}

The lemmas of the examples are axioms: the abstract predicates $\sorted$, $\perm$, $\partitioned$,
$\mathsf{bubbled}$, $\mathsf{shaken}$ and the functions $\argmin$, $\mathsf{insert}$, $\mathsf{shake}$,
$\mathsf{nodup}$, $\mathsf{merged}$, $\mathsf{lb}$ have no definitions, and the summaries state the
correctness of an inner loop in one line. Nothing in the tool checks them, so the discipline of the folder
replaces the check.
\begin{itemize}[nosep]
\item \textbf{Every result is checked on a wrong postcondition.} Each thread of \texttt{pending/} has a
  \texttt{program\_wrong.dfy} ($\sorted(b, 0, n + 1)$, the sum plus one, the negated flag); a prelude that
  closes its goal is unsound or vacuous, not a success. Two such cases occurred: a rule
  $\mathsf{nodup}(a, k, n) \Leftrightarrow \mathsf{true}$ under the hypothesis $\mathsf{distincts}(\ldots)$
  fired inside the context in which a side condition of a rule \emph{on} $\mathsf{distincts}$ was matched, so
  the atom was derived from itself (finding~13); and a condition with an open dependency and an equation on
  its bound variable was decided vacuously (\texttt{vac.dfy}). The lesson: a side condition must not be
  simplifiable by rules that use the conclusion, and the safe form of a condition is an atom that is matched
  literally.
\item \textbf{A lemma must be a lemma.} The lower bound closes with an exit $h \le l$ of two open variables
  matched against the fact $hi \le lo$ (\texttt{lb\_cmp0.dfy}); in this goal its conclusion is right, since
  the exit holds in the final state, but $h$ and $l$ are not tied to the traces and the rule is not a
  theorem. Such a rule is not admitted, and the file stays open (Section~\ref{sec:lb}).
\item \textbf{An induction constrains all steps.} A quantified condition over $0 \le k < n$ with
  $n \le \intf(m)$ steps leaves the further steps unconstrained. The number of steps is fixed exactly
  ($n = \intf(m)$, a fact of the goal), or the invariant is clamped so that every step keeps it
  ($\mathsf{bubbled}$ for $k \ge len$), or the result holds for any number of iterations beyond a bound
  ($m + n \le \intf(nn)$ in the merge).
\item \textbf{Soundness premises are the matcher's facts.} The Lean proofs of the core modules showed the
  premises that the lemmas need and the statements omit: dependencies of equal length, indices within the
  lists, complete guards; the fallback encoding does not assume them. The modules \texttt{nested},
  \texttt{sorting} and \texttt{strides} have no Lean proofs yet, and the abstract predicates would need
  definitions there for the step lemmas to be provable; both are open items.
\item \textbf{A killed run prints no goal.} A runner that counts goals reports a timeout as a proof: the
  first ``closed'' runs of \texttt{row\_sums}, \texttt{transpose} and most variants of the lower bound were
  kills at 120 or 300 s (finding~18). Every status of the folder was re-established with the exit status
  read, and long runs use a private copy of the jar, since a rebuild during a run crashes the JVM
  (finding~33).
\end{itemize}

\subsection{Termination}\label{sec:principles:termination}

Rewriting cycles by one mechanism, recorded for \lemma{sum\_invs\_const} of the database and met four
times in the folder (findings~19, 24, 29): an open $\mathsf{all}$ side condition is matched against the
$\mathsf{all}$ facts, which builds the simplifier's context of each fact's combinator and rewrites the reads
$\mathsf{earlier}(\natf(k), \bl{X})$ of its dependencies there; if the left-hand side of the rule being
decided matches one of those rewritten reads, the rule recurs through the same fact one level deeper until
the step budget is exhausted. Two ingredients, both needed: an $\mathsf{all}$ condition over lists that the
left-hand side binds, and a left-hand side general enough to match inside that context. Hence the rule of
thumb for a rule with $\mathsf{all}$ conditions: its left-hand side must be a term that no rewriting of a
read can produce, an $\mathsf{all}$ fragment with a specific body or an atom of the specification, never a
bare read, comparison, equation, or $\cur$ of a named trace for which an $\mathsf{earlier}$ acceleration
lemma exists or may be added. The general substitution lemmas of the first version of \texttt{row\_sums.dfy}
(a higher-order body $t = F(k, J)$, a read $\mathsf{earlier}(\natf(i), ts)$) are exactly what this forbids
(Section~\ref{sec:rows}); the exit lemma of the module \texttt{nested} with the body $n \le j$ cycled with
each counter lemma and with none alone, found by pair runs, and is stated on the constraint fragment
$\intf(nn) \le n$ now. The order of the modules also decides which abstraction a fragment gets:
\lemma{consts\_prefix\_def} of \texttt{aggregation} abstracts the exit $j = n$ before \texttt{nested} sees it,
hence \lemma{consts\_prefix\_integer} with the atom as left-hand side (finding~23).

\subsection{Diagnostics}\label{sec:principles:diag}

The obstacles of this folder were located by four techniques, each worth a flag of the tool:
\begin{itemize}[nosep]
\item \textbf{Preludes.} A program without lemmas and a file of lemmas spliced in front of it by the
  script \texttt{with.sh} of \texttt{pending/}, so that variants of the lemmas run against the same program;
  the folder holds the series for the cocktail sort (\texttt{cock2} to \texttt{cock9b}), the first duplicate
  (\texttt{fd4} to \texttt{fd8}), the lower bound, the row sums and the transpose, and the bisection of the
  module \texttt{nested}.
\item \textbf{Truncated copies.} A summary with its first $k$ side conditions, $k$ downwards, finds the first
  failing condition (the 21st of the cocktail summary); the copies are unsound and serve the diagnosis only.
\item \textbf{Probes.} Rules with the real left-hand side, the first $k$ conditions and a fresh atom
  $\mathsf{probe}_k(\ldots)$ as right-hand side show in the goal the largest $k$ whose conditions hold, and a
  probe that fires with ``variables of a rule not instantiated'' shows a left-hand side that left variables
  unbound (finding~32). A flag \texttt{-probe <lemma>} that reports the first failing condition of a named
  rule, with the facts it was matched against, would replace an afternoon of probe files.
\item \textbf{Pair runs.} The single lemmas of a module in pairs on the program that cycles find the pair
  that does (\texttt{pending/nested\_module/pairs.sh}); the rule reported at the step limit is the last step
  of the loop, not its cause.
\end{itemize}
With \texttt{-rules} a proof reports the lemmas it used; a residual goal points to the missing one. The
whole folder runs in about five minutes on one core, 5 to 35 s per file, the row sums and the transpose in
under two minutes.

\subsection{The same principles on C programs: the MonoCera tests}\label{sec:principles:monocera}

The 125 SV-COMP tests of MonoCera (Amilon et al., CAV 2023; instrumentation of $\forall$, $\max$, $\min$ and
$\sum$ over arrays) are C programs with a loop that fills an array and a loop that checks an assertion over
it. Their summary (\texttt{benchmarks/monocera/summary.md}) and the case studies of the lemmas written for
them (\texttt{lemmas.tex} there, Cuvée 0.35.15) show the principles of this folder on another class of
programs, without hand-written specifications.

\begin{table}[t]
\centering\small
\begin{tabular}{@{}lrrrrr@{}}
\toprule
operator & tests & MonoCera & Cuvée \texttt{-teq} & Cuvée \texttt{-chc} & TriCera, CPAchecker, SeaHorn \\
\midrule
forall & 93 & 26 & 53 & 0 & 0, 1, 0 \\
max & 4 & 2 & 2 & 0 & 0, 0, 0 \\
min & 11 & 3 & 7 & 0 & 0, 0, 0 \\
sum & 17 & 10 & 17 & 0 & 0, 0, 0 \\
\midrule
total & 125 & 41 & 79 & 0 & 0, 1, 0 \\
\bottomrule
\end{tabular}
\caption{Proved tests of the MonoCera evaluation (Cuvée 0.40.0, 15 min, Oct 12, 2026; the other tools' numbers
from the paper). All proofs of \texttt{-teq} end within 20 s (median 9 s); no tool gives a wrong result.}
\label{tab:monocera}
\end{table}

\begin{itemize}[nosep]
\item \textbf{The shape of the data decides the lemma, not the idea.} An accumulator kept in the cell
  \texttt{sum[0]} of a one-element array is a trace of stores, which \lemma{current\_trace\_reset\_add} of the
  database, written for a scalar, does not match; \lemma{select\_current\_trace\_reset\_add} is the same lemma
  for the read of the cell, with the scalar lemma's right-hand side so that the database continues from
  there. This is ``name the normal form'' (Section~\ref{sec:principles:shape}): the Dafny translations of
  SV-COMP make the cell a scalar and close the test for free, the C programs need the variant.
\item \textbf{Left-hand sides bind, side conditions must be facts.} \lemma{sum\_cases\_below} bounds a sum of
  $n$ terms that are each $0$, $c$ or $d$ by $e$ with its left-hand side a comparison whose constants bind
  the bound; the general \lemma{sum\_below\_prefix} of the database does not apply, since its bound occurs
  only in a side condition, which would have to be among the facts. The counting sort's
  \lemma{rank\_table\_def} binds the key bound $kk$ for the same reason (Section~\ref{sec:cs}), and the
  pattern $s = m + n$ of the merge is its arithmetic form.
\item \textbf{Lemmas compute values, they do not make goals vacuous.} Each closed test was rerun with a
  false assertion, which stays open with a residual such as $\neg(0 < N)$: the lemma computed $N - 1$ and the
  bound was wrong. The same check as the wrong postconditions of Section~\ref{sec:principles:sound}.
\item \textbf{Specific lemmas are summaries.} The strided stores (\lemma{select\_current\_trace\_stride2},
  stride 3, three exhaustive cases at stride 2) and the stores ahead ($F(k)$ at $k + 1$, then $G(k)$ at $k$)
  fix the number of stores and of cases on their left-hand sides, as the summaries of this folder fix the
  fragment structure of one algorithm; strides 3 to 5 with per-cell cases, writes by inner loops and
  recurrences on earlier entries remain open. Both want what the arity variants \lemma{perm\_steps},
  \lemma{perm\_steps\_two} and \lemma{sorted\_prefix\_steps\_two} want: a pattern with the stride, the number
  of dependencies, or the number of cases as a parameter.
\item \textbf{The checking loop is an instantiation at a witness.} In 72 of the 96 open goals of the
  earlier run the assertion of the checking loop survives as an $\mathsf{all}$ over its earlier iterations,
  and the goal is the assertion at the current index $i = \intf(n)$ of a filled array whose closed form is
  not derived (recurrences $b[i] = b[i - 1] + a[i - 1]$, parity, division): the step from the constraint of
  one loop to a position of another is the one that \lemma{not\_occurs\_selects\_def} and
  \lemma{select\_not\_occurs\_selects} make for the linear search (Section~\ref{sec:lsearch}), an atom and
  its instantiation, and the missing closed forms are the partial descriptions of Section~7 of the paper
  (the $\max$ and $\min$ tests, the lower bound here).
\item \textbf{The higher-order variable and the array.} The simplifier's eta reduction turned the instance
  $(z) \mapsto a[z]$ of a higher-order summand $G$ into the array $a$, leaving an application $a(k)$ that no
  rule about $a[k]$ matches; restricting eta reduction to lambdas of array type fixed it in 0.35.15. The
  examples of this folder meet the same phenomenon in the matcher (finding~10) and avoid a higher-order
  variable for a bare read, reading the array literally (\texttt{prefix\_sums.dfy}) or pinning the stored
  values (Section~\ref{sec:merge}).
\item \textbf{A match that fails for no visible reason is usually one of the known patterns.} The lemma for
  the stores ahead of a loop that runs downwards had the same printed form as the trace in the goal and did
  not match, and the summary of the MonoCera tests had excluded the literal $-1$ as the cause. Re-traced here
  (Section~\ref{sec:rewnif}), it is two of the patterns of Section~\ref{sec:principles:matcher} at once: the
  literal step $-1$ of the counter dependency is not matched (a variable $d$ with $d = -1$ as a side
  condition is, as in \texttt{reverse.dfy}), and the position $i - 1$ of the store ahead is arithmetic on a
  bound variable (the pin). With both, the lemma fires and the two \texttt{rewnifrev} tests close
  (\texttt{rewnifrev2} stores ahead upwards in a loop that runs downwards, another shape). The cocktail
  summary that fired once in a day (finding~12) and the lemma of the module \texttt{nested} that fires in
  the input but not from the module (finding~34) are residues of the same kind, for which the probe flag of
  Section~\ref{sec:principles:diag} is the remedy.
\item \textbf{Lookup against inference.} \texttt{-chc} with the Horn solver of z3 proves none of the 125
  tests (120 timeouts): the quantified invariants over arrays that the tests need, $b[k] = k^2 + k + 1$ for
  all $k < i$, are not found, as for TriCera in the paper, and the 15 minutes add nothing to either
  MonoCera or \texttt{-teq}, whose proofs end within 20 s. This is the paper's thesis in numbers: a lemma
  found once serves every loop of its shape, and the price is that a loop without a lemma stays open
  (43 of the 125, mostly the recurrences, the parity and the nondeterministic input arrays).
\end{itemize}

\subsection{What the engine and the database need}\label{sec:principles:next}

The workarounds of Section~\ref{sec:principles:matcher} are each a missing feature of the engine, ordered
by the payoff that the examples measure: a re-entrance guard, such that a rule is not applied while the
instances of its own side conditions are simplified (the sound lower bound closes, the summaries lose the
$\cur(\mathsf{results}(\ldots))$ detour); closed $\mathsf{all}$ conditions matched modulo unused
dependencies instead of decided (the fresh variables and identifying equations disappear from five files,
the summaries shrink by a third); a left-hand side matched through the definition of a trace variable binds
all its variables or is refused; a static check that rejects a left-hand side with a bare bound variable of a
variable dependency together with an $\mathsf{all}$ condition (it would have refused the cycling exit lemma
on load); and a case split of \texttt{-split} on a disjunctive negated exit. For the database: the module
\texttt{nested} into \texttt{-teq} with a run on the SV-COMP translations; the chain of the row sums
generalized into it (counters over constant numbers of iterations, final values given per iteration,
pointwise definitions as maps, reads of a map over the counter); the module \texttt{sorting}, which no file
includes yet, taken up by the examples; Lean proofs for \texttt{nested}, \texttt{sorting} and
\texttt{strides}, with definitions of the abstract predicates; and lemma schemas, one per observer of a
named scan, in place of the arity and type variants. The algorithms that remain open need, respectively,
the substitution of a solved exit into a function argument (\texttt{is\_sorted.dfy}), a witness for an
existential postcondition or quantified left-hand sides (\texttt{max\_element.dfy}), and the re-entrance
guard (the lower bound).
